% GENERATED FILE. Edit canonical lessons, then run npm run generate:books.
% canonical-source-hash: fnv1a64:7311898a3e84a046
% canonical-lessons: ES-C734-quemarse, ES-C734-quemadura, ES-C734-venda, ES-C734-sangrar, ES-C734-cicatriz

\chapter{Quemaduras y cortes}
\label{ch:es-quemaduras-y-cortes}
\hlchaptermodality{734}

\begin{chapteropening}
\textbf{By the end of this chapter:} \emph{I can describe a burn or a cut that bleeds, the bandage and the scar.}

Complete the last lesson of chapter 734: I can describe a burn or a cut that bleeds, the bandage and the scar.
\end{chapteropening}

\section[quemarse]{quemarse — to burn oneself; to get sunburnt}

\label{lesson:ES-C734-quemarse}



\begin{quote}
Before the new one: say the Spanish for to infect, then the Spanish for to sneeze.
\end{quote}

\subsection*{You'll want to know: quemarse}
\textbf{quemarse} — \textquotedblleft{}to burn oneself; to get sunburnt\textquotedblright{}.

\begin{quote}
Me quemé la mano con la sartén.
\end{quote}

\begin{quote}
\emph{I burned my hand on the frying pan.}
\end{quote}

\emph{Quemarse} + body part. \emph{Quemarse al sol} — to get sunburnt.

\begin{cousinweb}
\emph{quemarse} is the reflexive of \emph{quemar}, which is inherited from Latin \emph{cremāre} \textquotedblleft{}to burn\textquotedblright{} (EtymDB); English \emph{cremate} is the same Latin verb.
\end{cousinweb}

\begin{grammarlens}[title={What you've built}]
One more word for telling, explaining and describing at B1.
\end{grammarlens}

\subsection*{Your turn}
\begin{itemize}
\raggedright
  \item \emph{Say it:} \emph{quemarse}
  \item \emph{Say it:} \emph{quemarse}, once more
  \item \emph{Say it:} a sentence of your own with \emph{quemarse}
  \item \emph{Recall:} say the Spanish for to infect, then the Spanish for to sneeze, then say \emph{quemarse} again
\end{itemize}

\begin{tcolorbox}[breakable,colback=teal!4,colframe=teal!35!black,title={Before you move on}]
What does quemarse mean? (\textquotedblleft{}To burn oneself; to get sunburnt\textquotedblright{}.) Say it once more.
\end{tcolorbox}
\section[la quemadura]{la quemadura — the burn}

\label{lesson:ES-C734-quemadura}



\begin{quote}
Before the new one: say the Spanish for to sneeze, then the Spanish for to burn oneself.
\end{quote}

\subsection*{You'll want to know: la quemadura}
\textbf{la quemadura} — \textquotedblleft{}the burn\textquotedblright{}.

\begin{quote}
Tengo una quemadura en el brazo.
\end{quote}

\begin{quote}
\emph{I have a burn on my arm.}
\end{quote}

\begin{cousinweb}
\emph{quemadura} is built on \emph{quemar} \textquotedblleft{}to burn\textquotedblright{} with the ending \emph{-dura}; \emph{quemar} is inherited from Latin \emph{cremāre} (EtymDB).
\end{cousinweb}

\begin{grammarlens}[title={What you've built}]
One more word for telling, explaining and describing at B1.
\end{grammarlens}

\subsection*{Your turn}
\begin{itemize}
\raggedright
  \item \emph{Say it:} \emph{la quemadura}
  \item \emph{Say it:} \emph{la quemadura}, once more
  \item \emph{Say it:} a sentence of your own with \emph{la quemadura}
  \item \emph{Recall:} say the Spanish for to sneeze, then the Spanish for to burn oneself, then say \emph{la quemadura} again
\end{itemize}

\begin{tcolorbox}[breakable,colback=teal!4,colframe=teal!35!black,title={Before you move on}]
What does la quemadura mean? (\textquotedblleft{}The burn\textquotedblright{}.) Say it once more.
\end{tcolorbox}
\section[la venda]{la venda — the bandage}

\label{lesson:ES-C734-venda}



\begin{quote}
Before the new one: say the Spanish for to burn oneself, then the Spanish for the burn.
\end{quote}

\subsection*{You'll want to know: la venda}
\textbf{la venda} — \textquotedblleft{}the bandage\textquotedblright{}.

\begin{quote}
La enfermera me puso una venda.
\end{quote}

\begin{quote}
\emph{The nurse put a bandage on me.}
\end{quote}

A small sticking plaster is \emph{una tirita} (Spain) or \emph{una curita} (Latin America).

\begin{cousinweb}
\emph{venda} is usually traced to Germanic \emph{binda} \textquotedblleft{}band, strip of cloth\textquotedblright{}, from the root that gives English \emph{bind} (EtymDB gives the Indo-European root \emph{b\textsuperscript{h}end\textsuperscript{h}-} \textquotedblleft{}to bind\textquotedblright{}).
\end{cousinweb}

\begin{grammarlens}[title={What you've built}]
One more word for telling, explaining and describing at B1.
\end{grammarlens}

\subsection*{Your turn}
\begin{itemize}
\raggedright
  \item \emph{Say it:} \emph{la venda}
  \item \emph{Say it:} \emph{la venda}, once more
  \item \emph{Say it:} a sentence of your own with \emph{la venda}
  \item \emph{Recall:} say the Spanish for to burn oneself, then the Spanish for the burn, then say \emph{la venda} again
\end{itemize}

\begin{tcolorbox}[breakable,colback=teal!4,colframe=teal!35!black,title={Before you move on}]
What does la venda mean? (\textquotedblleft{}The bandage\textquotedblright{}.) Say it once more.
\end{tcolorbox}
\section[sangrar]{sangrar — to bleed}

\label{lesson:ES-C734-sangrar}



\begin{quote}
Before the new one: say the Spanish for the burn, then the Spanish for the bandage.
\end{quote}

\subsection*{You'll want to know: sangrar}
\textbf{sangrar} — \textquotedblleft{}to bleed\textquotedblright{}.

\begin{quote}
Me sangra la nariz.
\end{quote}

\begin{quote}
\emph{My nose is bleeding.}
\end{quote}

Works like \emph{doler}: \emph{me sangra la rodilla}.

\begin{cousinweb}
\emph{sangrar} is inherited from Latin \emph{sanguināre} \textquotedblleft{}to bleed\textquotedblright{} (EtymDB), built on \emph{sanguis} \textquotedblleft{}blood\textquotedblright{}.
\end{cousinweb}

\begin{grammarlens}[title={What you've built}]
One more word for telling, explaining and describing at B1.
\end{grammarlens}

\subsection*{Your turn}
\begin{itemize}
\raggedright
  \item \emph{Say it:} \emph{sangrar}
  \item \emph{Say it:} \emph{sangrar}, once more
  \item \emph{Say it:} a sentence of your own with \emph{sangrar}
  \item \emph{Recall:} say the Spanish for the burn, then the Spanish for the bandage, then say \emph{sangrar} again
\end{itemize}

\begin{tcolorbox}[breakable,colback=teal!4,colframe=teal!35!black,title={Before you move on}]
What does sangrar mean? (\textquotedblleft{}To bleed\textquotedblright{}.) Say it once more.
\end{tcolorbox}
\section[la cicatriz]{la cicatriz — the scar}

\label{lesson:ES-C734-cicatriz}



\begin{quote}
Before the new one: say the Spanish for the bandage, then the Spanish for to bleed.
\end{quote}

\subsection*{You'll want to know: la cicatriz}
\textbf{la cicatriz} — \textquotedblleft{}the scar\textquotedblright{}.

\begin{quote}
Tengo una cicatriz en la mano.
\end{quote}

\begin{quote}
\emph{I have a scar on my hand.}
\end{quote}

Plural \emph{las cicatrices}.

\begin{cousinweb}
\emph{cicatriz} comes from Latin \emph{cicātrīx}, \emph{cicātrīcem} \textquotedblleft{}scar\textquotedblright{} (EtymDB).
\end{cousinweb}

\begin{grammarlens}[title={What you've built}]
One more word for telling, explaining and describing at B1.
\end{grammarlens}

\subsection*{Your turn}
\begin{itemize}
\raggedright
  \item \emph{Say it:} \emph{la cicatriz}
  \item \emph{Say it:} \emph{la cicatriz}, once more
  \item \emph{Say it:} a sentence of your own with \emph{la cicatriz}
  \item \emph{Recall:} say the Spanish for the bandage, then the Spanish for to bleed, then say \emph{la cicatriz} again
\end{itemize}

\begin{tcolorbox}[breakable,colback=teal!4,colframe=teal!35!black,title={Before you move on}]
What does la cicatriz mean? (\textquotedblleft{}The scar\textquotedblright{}.) Say it once more.
\end{tcolorbox}
